\section{Project Health Monitoring}
\textit{Reference $|$ Course Text: Appendix B}

Detecting a struggling project early is the most valuable skill a PA can develop. By the time a project is visibly failing (missed design review, client complaint, team member dropout), the window for effective intervention has often closed. This section provides tools for ongoing health monitoring.

\subsection{The Project Health Check as a Management Tool}

The Course Text (Appendix~B) includes a Project Health Check: a self-assessment tool that asks teams to rate themselves Green/Yellow/Red across twelve common failure patterns. While designed for student self-reflection, the Health Check is equally valuable as a PA management tool.

\begin{paaction}
\textbf{Recommendation:} Make the Project Health Check a mandatory, graded deliverable at the end of each semester (Sprint~6 and Sprint~14). Require teams to complete it honestly and discuss the results with you. The act of self-diagnosis is as valuable as the resulting data.
\end{paaction}

\subsection{Green / Yellow / Red Indicators}

Use the following general framework to assess project health at any point in the semester:

\begin{faqbox}[Green: On Track]
\begin{itemize}[nosep,leftmargin=1.5em]
\item Deliverables are submitted on time and at or above the expected quality level.
\item Team meets regularly and all members participate actively.
\item Client communication is proactive and documented.
\item Schedule shows margin for testing and iteration.
\item The team can articulate their current risks and mitigation strategies.
\end{itemize}
\end{faqbox}

\begin{faqbox}[Yellow: Caution]
\begin{itemize}[nosep,leftmargin=1.5em]
\item One or two deliverables have been late or below expected quality.
\item Team meetings are inconsistent, or one member's attendance is spotty.
\item Client communication is reactive (responding to client inquiries rather than initiating).
\item Schedule shows no margin; the critical path has no float.
\item The team acknowledges risks but has no concrete mitigation plan.
\end{itemize}
\end{faqbox}

\begin{faqbox}[Red: Intervention Required]
\begin{itemize}[nosep,leftmargin=1.5em]
\item Multiple deliverables are late or substantially below quality expectations.
\item Team conflict is unresolved and affecting productivity.
\item Client has expressed concern or is unresponsive due to loss of confidence.
\item Schedule shows the team cannot complete the project as scoped.
\item The team cannot articulate what they are doing next or why.
\end{itemize}
\end{faqbox}

\subsection{When to Intervene vs.\ When to Let the Team Struggle}

Not all struggle is bad. The distinction between productive and destructive struggle is central to effective PA mentoring.

\textbf{Productive struggle} (let it continue):
\begin{itemize}[nosep,leftmargin=1.5em]
\item The team is wrestling with a genuine technical challenge and making incremental progress.
\item Team members are debating design trade-offs with data and reasoning (even if the debate is heated).
\item The team has identified a problem and is working through a solution, even if slowly.
\item Failure of a PoC or test provides learning that redirects the design.
\end{itemize}

\textbf{Destructive struggle} (intervene):
\begin{itemize}[nosep,leftmargin=1.5em]
\item The team is stuck on the same problem for more than one sprint with no visible progress.
\item Conflict has become personal, and team members are avoiding each other rather than collaborating.
\item The team does not recognize they are in trouble (no self-awareness of project health).
\item Safety is at risk: fabrication without a Hazard Assessment, skipping test procedures, or ignoring equipment limitations.
\item The team has given up on a deliverable or a requirement without discussing it with you.
\end{itemize}

\begin{bestpractice}
When you decide to let a team struggle productively, tell them: ``I can see this is challenging, and I think working through it will teach you more than if I give you the answer. I am here if you get completely stuck.'' This acknowledges the difficulty without removing agency. Check in at the next meeting to verify progress.
\end{bestpractice}

\subsection{Documentation Debt: The Leading Indicator}

If you track only one health metric, track documentation debt.

\begin{redflag}
Documentation debt is the number one indicator of a project in trouble. Teams that build without writing---deferring all report sections, test logs, and design rationale to ``later''---produce rushed, low-quality final deliverables and often miss critical details because they are reconstructing decisions from memory rather than from contemporaneous records.
\end{redflag}

\textbf{How to monitor:}
\begin{itemize}[nosep,leftmargin=1.5em]
\item At each sprint review, ask: ``What report sections have you written this sprint?''
\item Check the team's shared document (Google Docs, Overleaf, or similar) for recent activity.
\item Compare the number of pages drafted against the number of pages expected at this point in the semester.
\item If the team is more than two sprints behind on documentation, flag it as Yellow and establish a remediation plan.
\end{itemize}

\begin{paaction}
Establish a documentation expectation in Sprint~0: report sections should be drafted during the sprint in which the work is performed, not retroactively at the end of the semester. Reinforce this expectation at every sprint review.
\end{paaction}


% ═══════════════════════════════════════
